\documentclass{article}
\usepackage{amsmath,amssymb,amsthm}
\usepackage[margin=1in]{geometry}
\usepackage{array}

\newtheorem{theorem}{Theorem}
\newtheorem{conjecture}{Conjecture}
\newtheorem{claim}{Claim}

\title{Disproof of Conjecture 00000000423:\\Promotion Orbit Lengths Do Not Always Divide the Tableau Size}
\author{TLMC submission}
\date{October 2, 2026}

\begin{document}
\maketitle

\begin{abstract}
Conjecture 00000000423 asserts that promotion orbit lengths on standard Young
tableaux divide the tableau size. We disprove it with the shape $(2,1)$: the
two standard Young tableaux of that shape form a single promotion orbit of
length $2$, while the tableau size is $n=3$ cells, and $2\nmid 3$. The
computation is three lines of jeu-de-taquin; it is re-verified by an
independent script (\texttt{reproduce.py}) and formalized in core Lean~4 with
zero axioms (\texttt{lean4/}).
\end{abstract}

\section{The conjecture}

For a partition $\lambda$ of $n$, let $\mathrm{SYT}(\lambda)$ denote the set
of standard Young tableaux of shape $\lambda$. Sch\"utzenberger's
\emph{promotion} operator $\pi$ acts on $\mathrm{SYT}(\lambda)$ by cyclic
sliding: delete the entry $1$, perform jeu-de-taquin slides into the vacated
cell, subtract $1$ from every remaining entry, and place $n$ in the (corner)
cell vacated by the slide.

\begin{conjecture}[00000000423]
Promotion orbit lengths divide the tableau size; the simultaneous
equidistribution and orbit counting are given jointly by a bijection.
\end{conjecture}

The first clause is the target of this disproof: orbit lengths of $\pi$ on
$\mathrm{SYT}(\lambda)$ should divide the tableau size.

\section{The attack: shape $(2,1)$}

Shape $(2,1)$ has $n=3$ cells and exactly two standard fillings
(entries $1,2,3$; rows and the first column strictly increasing):
\[
A=\begin{array}{|c|c|}\hline 1&2\\\cline{1-2}\multicolumn{1}{|c|}{3}&\multicolumn{1}{c}{}\\\cline{1-1}
\end{array}
\qquad\qquad
B=\begin{array}{|c|c|}\hline 1&3\\\cline{1-2}\multicolumn{1}{|c|}{2}&\multicolumn{1}{c}{}\\\cline{1-1}
\end{array}
\]

\begin{claim}[$\pi(A)=B$]\label{cl:A}
Delete the $1$ from $A$; the empty cell $(1,1)$ has right neighbour $2$ and
lower neighbour $3$; the smaller, $2$, slides left. The vacated cell $(1,2)$
has no right or lower neighbour, so the slide stops. Subtract $1$ from the
remaining entries ($3\mapsto 2$) and place $3$ at $(1,2)$:
\[
A\;\longmapsto\;\begin{array}{|c|c|}\hline 2&\cdot\\\cline{1-2}\multicolumn{1}{|c|}{3}&\multicolumn{1}{c}{}\\\cline{1-1}\end{array}
\;\longmapsto\;\begin{array}{|c|c|}\hline 1&3\\\cline{1-2}\multicolumn{1}{|c|}{2}&\multicolumn{1}{c}{}\\\cline{1-1}\end{array}=B.
\]
\end{claim}

\begin{claim}[$\pi(B)=A$]\label{cl:B}
Delete the $1$ from $B$; the empty cell has right neighbour $3$ and lower
neighbour $2$; the smaller, $2$, slides up. The vacated cell $(2,1)$ has no
neighbour, so the slide stops. Subtract $1$ ($3\mapsto 2$) and place $3$ at
$(2,1)$, giving $A$.
\end{claim}

Hence $\pi$ is the transposition $A\leftrightarrow B$ on
$\mathrm{SYT}(2,1)$: the entire tableau set is a \emph{single orbit of length
$2$} (and $\pi\neq\mathrm{id}$ since $A\neq B$).

\begin{theorem}[Disproof of 00000000423]\label{thm:main}
$\pi$ on $\mathrm{SYT}(2,1)$ has orbit length $2$, the tableau size is $n=3$,
and $2\nmid 3$ ($3 = 2\cdot 1 + 1$). Therefore the conjecture is false.
\end{theorem}

\section{Boundary of the attack}

\begin{itemize}
\item \textbf{Rectangles are not counterexamples.} For an $a\times b$
rectangle with $n=ab$ cells, Haiman's formula gives promotion order
$n/\gcd(n,a)$, which divides $n$. The failure genuinely requires a
non-rectangular shape.
\item \textbf{Minimality.} Under the reading ``tableau size $=$ number of
cells,'' $(2,1)$ is the smallest counterexample: it is the smallest partition
that is not a row, column or rectangle with the divisibility holding.
\item \textbf{The alternative reading ``size $=$ \#SYT'' also fails.} For
shape $(3,2)$: $n=5$ cells, $\#\mathrm{SYT}=5$, and the promotion orbits have
lengths $\{3,2\}$ (lcm $6$). Neither $3$ nor $2$ divides $5$, and $6$ divides
neither $3$ nor $5$. So both natural readings of ``tableau size'' are
refuted, at $(2,1)$ and $(3,2)$ respectively.
\end{itemize}

\section{Verification}

The tableau enumeration and promotion computation were performed three
independent ways:
\begin{enumerate}
\item the jeu-de-taquin computation of Claims~\ref{cl:A} and~\ref{cl:B}
above;
\item the script \texttt{reproduce.py}, which enumerates $\mathrm{SYT}$ of a
shape by brute force and computes promotion orbits (shape $(2,1)$: orbit
lengths $[2]$; shape $(3,2)$: orbit lengths $[2,3]$, lcm $6$);
\item a core Lean~4 formalization (\texttt{lean4/}), whose main theorem
\texttt{conjecture\_00000000423\_false} states that the tableau set equals
$\{(1,2,3),(1,3,2)\}$, that promotion swaps the two, and that
$(3:\mathbb{N})\bmod 2 = 1$. \texttt{\#print axioms} on every theorem
reports ``does not depend on any axioms''; there are no \texttt{sorry}.
\end{enumerate}

\end{document}
